\chapter*{Vorwort}
Dieses Buch entstand aus einer einfachen Überzeugung: Wer einen Roboter programmieren will, muss Python nicht vollständig kennen -- er muss die richtigen Teile kennen, zur richtigen Zeit.

Der rote Faden: Ein Differentialantrieb-Roboter, der mit jedem Kapitel smarter wird. Am Ende des Buches fährt er autonom -- gesteuert von Python und ROS2.


\chapter*{Hinweise zur Benutzung}

\textbf{Code-Listings} zeigen Python-Code mit Zeilennummern. Alle Beispiele sind lauffähig.

\textbf{Farbige Boxen} signalisieren:
\begin{itemize}
    \item \textcolor{blue!70!black}{\textbf{Merksatz}} -- das Wichtigste auf einen Blick
    \item \textcolor{green!60!black}{\textbf{Praxisbezug}} -- direkte Robotik-Anwendung
    \item \textcolor{orange!80!black}{\textbf{Übung}} -- selbst ausprobieren
    \item \textcolor{purple!60!black}{\textbf{Analogie}} -- Vergleich aus dem Alltag
    \item \textcolor{red!70!black}{\textbf{ROS2-Verbindung}} -- Brücke zu ROS2
\end{itemize}

\textbf{Voraussetzungen:} Grundrechenarten, Interesse an Technik. Keine Programmiererfahrung nötig.
